\chapter{Conclusións e traballo futuro}
\label{chap:conclusions}

\lettrine{D}{erradeiro} capítulo da memoria, onde se presentará a situación final do traballo, as leccións aprendidas, a relación coas competencias da titulación en xeral e a mención en particular, posibles liñas futuras.

\section{Conclusións}
Reproducir música de forma colaborativa e equitativamente de forma dixital é posible. O obxectivo deste traballo foi a análise, deseño e implementación dunha gramola virtual. Para isto establecéronse uns requisitos funcionais, crear sesións colaborativas, amosar unha cola de reprodución, permitir engadir cancións a cola e reproducir a música. Uns requisitos de deseño, minimalista, fácil de comprender e fácil de usar. Estudáronse as diferentes opcións existentes, as carencias destas, as súas fortalezas e deseñouse unha plataforma que permitise a reprodución de música de forma colaborativa en distintas sesións na que os usuarios poden engadir novas cancións e votar polas súas favoritas.\smallbreak
Finalmente, comprobouse que o sistema creado cumpre os requisitos desexados, creouse unha gramola virtual que permite aos usuarios o uso mais básico de estas, que é a reprodución das cancións preferidas nunha orde, engadir cancións para a súa posterior reprodución e a posibilidade de votar para cambiar a orde de reprodución. Ademais de ser un sistema o mais sinxelo posible o que facilita o uso, con poucas iteracións o usuario pode crear unha sesión e reproducir música en calquera dispositivo con Spotify de forma colaborativa.

\section{Traballo futuro}
Aínda que se obtivo un produto completamente funcional, existe a posibilidade da implementación de novas ideas que melloren o produto final, neste caso a lista podería ser moi longa xa que existen moitos sistemas diferentes para a reprodución de música e a inclusión de algúns destes melloraría o valor da aplicación. Nun principio a aplicación non busca o lucro, aínda que se poderían poñer en funcionamento algunhas funcionalidades para que así fose. A continuación expóñense algunhas das ideas que se poderían desenvolver para a mellora da aplicación, estas son sempre opcións non exclusivas, polo que cun panel de axustes o administrador podería elixir a súa configuración desexada. 
\begin{itemize}
	\item \textbf{Modo DJ:} cun funcionamento similar, a aplicación non engadiría as cancións automaticamente a cola, mostraría unha lista a unha persoa que se dedique a pinchar segundo os votos dos usuarios e as súa propia elección.
	\item \textbf{Sistema de créditos:} limitar o número de cancións por usuario e ofrecer un código para conseguir aumentar o límite, por exemplo nun negocio de restauración para incitar o consumo dos clientes.
	\item \textbf{Usuarios globais:} poder ter o mesmo usuario para diferentes sesións en distintos lugares. Deste xeito poderíase manter un avatar e unhas estatísticas de uso persoais.
	\item \textbf{Estatísticas:} amosar aos usuarios estatísticas da sesión, cancións mais reproducidas, usuarios que mais cancións engaden\dots
	\item \textbf{Gardar listas:} as cancións que soen nun período de tempo gardalas nunha lista de reprodución de Spotify, deste xeito poderíase gardar por exemplo a lista de reprodución dunha noite nun bar.
	\item \textbf{Escoller estilos musicais:} permitir ao administrador limitar os estilos musicais que devolven as procuras.
	\item \textbf{OAuth:} acceder con Facebook, Google...
	\item \textbf{Aplicación iOS, Android:} posibilidade de compartir directamente dende Spotify un enlace a unha canción.
\end{itemize}

